\label{sec:intro}

Mobile and multihomed hosts often see several access networks at once (for example Wi-Fi and cellular). Regular TCP binds a connection to a single four-tuple, so a change of interface or loss of that path can interrupt the application unless higher layers recover. Multipath TCP (MPTCP) extends TCP so one logical connection may use several \emph{subflows} across different paths, with the goal of higher throughput when paths add capacity and smoother handover when one path disappears~\cite{rfc8684,Bonaventure2012usecases}.

This term project builds a small \emph{experimental framework} to compare baseline single-path TCP against a multipath-style setup under controlled link conditions. We use Mininet~\cite{Lantz2010mininet} with \texttt{TCLink} bandwidth and delay, Linux policy routing and kernel MPTCP endpoints where applicable, and \texttt{iperf3}~\cite{iperf3} with JSON output for throughput time series. Mobility-like events are \emph{emulated}: sudden interface down on one path, bandwidth collapse via \texttt{tc}, and RTT inflation via \texttt{netem}~\cite{linuxnetem}.

\textbf{Objectives} (aligned with the course project brief):
\begin{itemize}[noitemsep]
  \item Quantify per-path and combined throughput under heterogeneous caps and delays.
  \item Observe session continuity when one path is removed while traffic continues on another.
  \item Stress one path with bandwidth collapse and RTT spikes and compare aggregate behavior to single-path TCP.
  \item Discuss latency and scheduler effects qualitatively from throughput traces.
  \item Study a handover-like scenario by ramping RTT on the preferred path while a second path stays available (Experiment~6).
\end{itemize}
\textbf{Competing-flow fairness} in the sense of multiple independent flows sharing the same paths (e.g., Jain fairness with cross-traffic) is \emph{not} measured here; Experiment~6 instead reports scheduler and throughput behavior under primary-path degradation.

\textbf{Contributions of this report:} (i)~a reproducible Mininet-based workflow mapping six Python experiment drivers to measurement outputs; (ii)~numeric summaries from our captured logs (including tabular time-slice data where no figure is available); and (iii)~an honest discussion of limitations, including that our multipath measurements use parallel \texttt{iperf3} flows bound to interfaces as a practical stand-in for a single MPTCP socket in all scenarios.
